\chapter{Thermoelastic laser source}\label{ch:source}
Let absorbed volumetric power be
\begin{equation}
 Q(x,z,t)=\eta_{\rm abs}P(t)g_w(x)\,a_\delta(z),\quad
 g_w=\frac{\ee^{-x^2/w^2}}{\sqrt\pi w},\quad
 a_\delta=\frac{\ee^{-(h-z)/\delta}}{\delta[1-\ee^{-2h/\delta}]},
 \label{eq:physical-Q}
\end{equation}
where $P$ is incident line power (\si{W.m^{-1}}), $\eta_{\rm abs}=1-R$ is
absorbed fraction, $w$ is the stated $1/e$ intensity half-width, and $\delta$ is
absorption depth. Both profiles integrate to unity, so \cref{eq:physical-Q} has
units \si{W.m^{-3}}. \TODO{Supply wavelength-dependent $R$, $\delta$, and the
measured beam convention.}

With one-way coupling, temperature satisfies
\begin{equation}
 \rho C_p\partial_t\Delta T-K_T(\partial_x^2+\partial_z^2)\Delta T=Q.
 \label{eq:heat}
\end{equation}
We initially impose thermally insulating faces,
$\partial_z\Delta T(\pm h,t)=0$, appropriate only when heat loss during the
acoustic record is negligible. \TODO{Replace by measured convection/contact
conditions when the experiment is fixed.} Fourier transformation in $x$ and
expansion in Neumann thickness eigenfunctions $\chi_m$ gives
\begin{equation}
 \dot T_m+\kappa(k^2+\beta_m^2)T_m=Q_m/(\rho C_p),\quad
 T_m(k,t)=\frac1{\rho C_p}\int_0^t\!Q_m(k,t')
 \ee^{-\kappa(k^2+\beta_m^2)(t-t')}\dd t',\label{eq:heat-solution}
\end{equation}
where $\kappa=K_T/(\rho C_p)$ and
$Q_m=\int_{-h}^h\chi_mQ\dd z$. Negligible diffusion during a short pulse means
$\kappa(k^2+\beta_m^2)\tau\ll1$ for all materially excited components, not merely
``a nanosecond pulse.''

Thermal strain is $\epsilon_{ij}^{\rm th}=\alpha\Delta T\delta_{ij}$ and
\begin{equation}
 \sigma_{ij}=C_{ijkl}(\epsilon_{kl}-\epsilon_{kl}^{\rm th})
 =\sigma_{ij}^{\rm el}-\underbrace{(3\lambda+2\mu)\alpha\Delta T\delta_{ij}}_{\tau_{ij}}.
 \label{eq:thermal-stress}
\end{equation}
Writing $\rho\ddot u_i-\sigma_{ij,j}^{\rm el}=-\tau_{ij,j}$ produces a body
force but also requires distributional surface terms if $\tau$ terminates at a
boundary. The safer weak form follows after multiplying by a test field $V_i^*$:
\begin{equation}
 \int\rho V_i^*\ddot u_i\dd z+\int V_{i,j}^*C_{ijkl}u_{k,l}\dd z
 =\int V_{i,j}^*\tau_{ij}\dd z,
 \label{eq:weak-thermal}
\end{equation}
when total traction is zero. The integration by parts has included the surface
term correctly.

For mass-normalized $\vect U_n(z;k)$, modal projection gives the physical force
\begin{equation}
 F_n(k,t)=\int_{-h}^{h}U_{n,i,j}^*(z;k)\tau_{ij}(k,z,t)\dd z,
 \quad [F_n]=\si{m.s^{-2}}
 \label{eq:physical-overlap}
\end{equation}
under normalization \cref{eq:modal-mass}. The $x$ derivative contributes
$-\ii kU_{n,x}^*$ and the $z$ derivative $\partial_zU_{n,z}^*$. Source parity
then controls S/A coupling. Top-surface absorption breaks mid-plane symmetry and
generally excites both families; symmetric two-sided heating can suppress the
parity-incompatible overlap. Beam width sets the $k$ spectrum, while absorption
depth weights modal dilatation through thickness. Appendix~\ref{app:projection}
states the equivalent body-force and weak forms side by side.
